% Abstract of the whole thesis. Front matter: after the acknowledgements,
% immediately before Contributions at a glance.
% Word count of the abstract body: 346 words.
\chapter*{\centering Abstract}
\phantomsection
\label{sup:abstract}
\addcontentsline{toc}{chapter}{Abstract}
\markboth{ABSTRACT}{ABSTRACT}

The early course of an infection can be decided by the fate of a few
pathogens, for which mean population growth alone does not determine
establishment. Motivated by the intracellular phase of
\textit{Yersinia pestis} infection, this thesis develops stochastic
branching models to investigate replication, host-cell rupture and the
consequences of releasing pathogens in bursts rather than through
continuous budding.

The central model is a birth--death--catastrophe process in which
individual pathogens replicate, die or trigger rupture of their shared
host cell. Generating-function and Riccati methods yield distributions
and moments of intracellular load, released numbers and rupture timing.
The burst-size distribution, conditional on rupture, is shown to
coincide with the geometric quasi-stationary distribution of the
intracellular population. A two-type extension with irreversible
conversion and type-specific catastrophe rates yields exact
hypergeometric expressions for the finite-time no-catastrophe
probability and the corresponding bivariate generating function.

Embedding the single-cell laws in an age-structured renewal system
gives exact population means under linear infection and clearance. The
standard within-host model emerges as a projection valid within an
exponential growth phase: its effective rates depend on that phase and
cannot uniquely identify the three intracellular rates. At matched mean
lifetime output, bursting and budding have the same basic reproduction
number, but can differ in establishment probability. In the
supercritical regime, bursting reduces extinction risk precisely when
$\delta(\lambda+\mu)>\lambda\mu$, where $\lambda$, $\mu$ and $\delta$
are the per-pathogen replication, death and rupture rates. This
ordering reverses for expected output surviving a fixed positive
integer removal budget applied to one delivery. Models of depletable
intracellular defences further quantify the survival advantage of
larger founding cohorts.

Complementary analyses identify the survival constant of subcritical
binary branching with a value of the logistic map's Koenigs function
and establish that function's hypertranscendence in its spatial
argument. Models with changing evolutionary rates yield a limiting
species-count distribution under capacity-dependent speciation and an
explicit moving critical point for competitive reversal. Applied to
\yp{}, the infection models yield hypotheses for the transition from
intracellular replication to resistance to phagocytosis. These remain
untested against plague data and require measurements of intracellular
growth, death and release.
